
\newpage
%%\addcontentsline{toc}{chapter}{\MakeUppercase{Results and conclusions}}
\thispagestyle{empty}
%
\begin{spacing}{1.5}
%


\end{spacing}
\newpage

\thispagestyle{empty}
%
\begin{spacing}{1.2}
%

\chapter{Performance and Security}

\section{Performance}

The Smart GECI system is designed to ensure efficient spatial data processing,
responsive map visualization, and reliable concurrent access within an institutional
environment.

\subsection*{Map Response Time}

\begin{itemize}
    \item The system maintains an average map response time of \textbf{3–5 seconds}
    for loading spatial layers and retrieving metadata.
    \item Spatial queries from PostGIS are optimized using GiST indexing.
    \item GeoServer caching mechanisms are utilized to reduce repeated rendering
    overhead.
\end{itemize}

This ensures smooth interaction during zooming, panning, and feature selection.

\subsection*{Concurrent User Support}

\begin{itemize}
    \item The system supports up to \textbf{100 concurrent users} without noticeable
    performance degradation.
    \item Backend API requests are handled asynchronously using Node.js’s
    event-driven architecture.
    \item Database query optimization minimizes latency under multi-user access.
\end{itemize}

This capacity is suitable for institutional-level deployment within campus networks.

\subsection*{Dataset Update Propagation}

\begin{itemize}
    \item Administrative dataset updates are reflected within \textbf{10 seconds}
    across the frontend interface.
    \item GeoServer layer refresh mechanisms ensure updated spatial data is
    republished promptly.
    \item Metadata changes are immediately accessible through API endpoints.
\end{itemize}

This near real-time update capability ensures data consistency and reliability.

\subsection*{Intelligent Query Performance}

\begin{itemize}
    \item Average response time for natural language queries is \textbf{2–6 seconds},
    depending on query complexity.
    
    \item Retrieval-based responses (RAG) execute faster due to vector similarity
    indexing.
    
    \item Structured NL2SQL queries are validated before execution to prevent
    inefficient full-table scans.
    
    \item Spatial query results are optimized by limiting geometry payload size
    before GeoJSON conversion.
\end{itemize}

The hybrid query architecture ensures AI-assisted querying without
significantly affecting core map performance.

\subsection*{Raster Analytics Performance (UHI Module)}

\begin{itemize}
    \item Raster layers are served as tiled map services to improve loading speed.
    
    \item Pixel-level value retrieval is optimized using preprocessed raster layers.
    
    \item Zonal statistics computation is performed asynchronously to prevent UI blocking.
    
    \item Regional-scale implementation under Idukki Taluk ensures meaningful
    resolution without excessive pixel-level overhead.
\end{itemize}

These optimizations allow environmental analytics to coexist with
infrastructure mapping without degrading system responsiveness.

\section{Security}

Security is a critical component of Smart GECI to protect institutional data,
prevent unauthorized access, and ensure accountability.

\subsection*{Secure Communication}

\begin{itemize}
    \item All client-server communication is conducted over \textbf{HTTPS}.
    \item TLS/SSL encryption protects data in transit between frontend,
    backend, and database layers.
\end{itemize}

\subsection*{Role-Based Authentication}

\begin{itemize}
    \item Role-Based Access Control (RBAC) differentiates permissions among:
        \begin{itemize}
            \item General Users (view-only access),
            \item Registered Users (authenticated advanced access),
            \item Administrators (full system control).
        \end{itemize}
    \item Administrative operations require authenticated access.
    \item Sensitive endpoints are protected using authentication middleware.
\end{itemize}

\subsection*{Access Logging and Update History}

\begin{itemize}
    \item All administrative actions are logged for accountability.
    \item Dataset updates are recorded with timestamps.
    \item Unauthorized access attempts are logged and monitored.
\end{itemize}

\subsection*{AI Query Validation and Safety}

\begin{itemize}
    \item All generated SQL queries are validated before execution to prevent
    injection or unsafe operations.
    
    \item Query execution is restricted to read-only operations for non-administrative users.
    
    \item Natural language inputs are sanitized before processing.
    
    \item Execution logs are maintained for intelligent query auditing.
\end{itemize}
